\subsection{tracelb: descubrimiento adicional y patrones}\label{sec:lb}
\subsubsection{Valparaíso}
\begingroup
\fontsize{9}{11}\selectfont
\addtocounter{table}{-1}
\begin{longtable}{p{3.2cm}p{8cm}p{3cm}}
\toprule
\textbf{Modo} & \textbf{IP adicionales r1 / r2 / r3} & \textbf{Mediana sondas} \\ \midrule
\endfirsthead
\toprule
\textbf{Modo} & \textbf{IP adicionales r1 / r2 / r3} & \textbf{Mediana sondas} \\ \midrule
\endhead
udp-dport & 0 / 0 / 0 & 98 \\
icmp-echo & 0 / 0 / 0 & 44 \\
udp-sport & 0 / 0 / 0 & 38 \\
tcp-sport & 0 / 0 / 0 & 38 \\
tcp-ack-sport & 0 / 0 / 0 & 27 \\
\bottomrule
\end{longtable}
\endgroup
La referencia es la unión de los seis trace de la misma ronda y destino. IP adicionales observadas: udp-dport: ninguna; icmp-echo: ninguna; udp-sport: ninguna; tcp-sport: ninguna; tcp-ack-sport: ninguna. Un aumento de IP revela interfaces adicionales bajo mayor exploración; no acredita por sí solo un menor costo por descubrimiento.

\subsubsection{La Habana}
\begingroup
\fontsize{9}{11}\selectfont
\addtocounter{table}{-1}
\begin{longtable}{p{3.2cm}p{8cm}p{3cm}}
\toprule
\textbf{Modo} & \textbf{IP adicionales r1 / r2 / r3} & \textbf{Mediana sondas} \\ \midrule
\endfirsthead
\toprule
\textbf{Modo} & \textbf{IP adicionales r1 / r2 / r3} & \textbf{Mediana sondas} \\ \midrule
\endhead
udp-dport & 2 / 2 / 0 & 214 \\
icmp-echo & 0 / 0 / 0 & 118 \\
udp-sport & 0 / 0 / 0 & 38 \\
tcp-sport & 0 / 0 / 0 & 38 \\
tcp-ack-sport & 0 / 0 / 0 & 27 \\
\bottomrule
\end{longtable}
\endgroup
La referencia es la unión de los seis trace de la misma ronda y destino. IP adicionales observadas: udp-dport: 154.54.168.85, 154.54.3.29, 69.79.104.1; icmp-echo: ninguna; udp-sport: ninguna; tcp-sport: ninguna; tcp-ack-sport: ninguna. Un aumento de IP revela interfaces adicionales bajo mayor exploración; no acredita por sí solo un menor costo por descubrimiento.

\subsubsection{Rusia}
\begingroup
\fontsize{9}{11}\selectfont
\addtocounter{table}{-1}
\begin{longtable}{p{3.2cm}p{8cm}p{3cm}}
\toprule
\textbf{Modo} & \textbf{IP adicionales r1 / r2 / r3} & \textbf{Mediana sondas} \\ \midrule
\endfirsthead
\toprule
\textbf{Modo} & \textbf{IP adicionales r1 / r2 / r3} & \textbf{Mediana sondas} \\ \midrule
\endhead
udp-dport & 1 / 0 / 0 & 204 \\
icmp-echo & 0 / 0 / 0 & 104 \\
udp-sport & 0 / 0 / 0 & 38 \\
tcp-sport & 0 / 0 / 0 & 38 \\
tcp-ack-sport & 0 / 0 / 0 & 27 \\
\bottomrule
\end{longtable}
\endgroup
La referencia es la unión de los seis trace de la misma ronda y destino. IP adicionales observadas: udp-dport: 195.22.211.209; icmp-echo: ninguna; udp-sport: ninguna; tcp-sport: ninguna; tcp-ack-sport: ninguna. Un aumento de IP revela interfaces adicionales bajo mayor exploración; no acredita por sí solo un menor costo por descubrimiento.

\Needspace{9\baselineskip}
\subsubsection{scamper / tracelb / udp-dport}
Ejemplo: Valparaíso. Entrenamiento: scamper\_\allowbreak{}20261026T204825\_\allowbreak{}aa1fe10f; contraste: rondas 2 y 3. Las reglas completas de los tres destinos están en el archivo de predicciones enlazado en el anexo.

\begingroup
\fontsize{9}{11}\selectfont
\addtocounter{table}{-1}
\begin{longtable}{p{3.1cm}p{5.8cm}p{5.1cm}}
\toprule
\textbf{Campo} & \textbf{Regla en ronda 1} & \textbf{Contraste posterior} \\ \midrule
\endfirsthead
\toprule
\textbf{Campo} & \textbf{Regla en ronda 1} & \textbf{Contraste posterior} \\ \midrule
\endhead
ip.ttl & variable en la captura & r2: no exacto; r3: no exacto \\
ip.id & incremento 1 por paquete & r2: cumple; r3: cumple \\
ip.len & constante 44 & r2: cumple; r3: cumple \\
ip.proto & constante 17 & r2: cumple; r3: cumple \\
udp.srcport & constante 32769 & r2: cumple; r3: cumple \\
udp.dstport & variable en la captura & r2: no exacto; r3: no exacto \\
udp.length & constante 24 & r2: cumple; r3: cumple \\
udp.checksum & incremento 1 por paquete & r2: cumple; r3: cumple \\
\bottomrule
\end{longtable}
\endgroup
Payload: [16] bytes. Hexadecimal inicial (offset 0):

\begin{quote}\ttfamily\footnotesize\raggedright 66b1 0000 0000 0000 0000 0000 0000 0000\end{quote}
La suma de payload\_\allowbreak{}word0\_\allowbreak{}BE, udp.checksum, udp.srcport, udp.dstport módulo 65535 se mantuvo constante. La base se aprende del primer paquete de cada ejecución; no se presume universal. Es compatible con compensación de checksum, sin probar la intención del emisor.

Predicción concreta en ronda 2, campo ip.id: la primera solicitud tiene valor 0x0001. Con la regla aprendida y esa base cuando corresponde, se predicen las solicitudes 2 a 4 antes de contrastarlas con sus valores observados.

\begingroup
\fontsize{9}{11}\selectfont
\addtocounter{table}{-1}
\begin{longtable}{p{3.2cm}p{5.1cm}p{5.1cm}}
\toprule
\textbf{Solicitud / trama} & \textbf{Predicción} & \textbf{Observación} \\ \midrule
\endfirsthead
\toprule
\textbf{Solicitud / trama} & \textbf{Predicción} & \textbf{Observación} \\ \midrule
\endhead
2 / 3 & 2 & 0x0002 \\
3 / 5 & 3 & 0x0003 \\
4 / 7 & 4 & 0x0004 \\
\bottomrule
\end{longtable}
\endgroup
Alcance: 93 solicitudes en el ejemplo inicial. Una verificación de campo o byte no equivale a una medición independiente; los campos variables conservan sus restricciones y no reciben valores inventados.

\Needspace{9\baselineskip}
\subsubsection{scamper / tracelb / icmp-echo}
Ejemplo: Valparaíso. Entrenamiento: scamper\_\allowbreak{}20261026T204853\_\allowbreak{}edf82eeb; contraste: rondas 2 y 3. Las reglas completas de los tres destinos están en el archivo de predicciones enlazado en el anexo.

\begingroup
\fontsize{9}{11}\selectfont
\addtocounter{table}{-1}
\begin{longtable}{p{3.1cm}p{5.8cm}p{5.1cm}}
\toprule
\textbf{Campo} & \textbf{Regla en ronda 1} & \textbf{Contraste posterior} \\ \midrule
\endfirsthead
\toprule
\textbf{Campo} & \textbf{Regla en ronda 1} & \textbf{Contraste posterior} \\ \midrule
\endhead
ip.ttl & variable en la captura & r2: no exacto; r3: no exacto \\
ip.id & constante 0x0000 & r2: cumple; r3: cumple \\
ip.len & constante 44 & r2: cumple; r3: cumple \\
ip.proto & constante 1 & r2: cumple; r3: cumple \\
icmp.type & constante 8 & r2: cumple; r3: cumple \\
icmp.code & constante 0 & r2: cumple; r3: cumple \\
icmp.ident & constante 32769 & r2: cumple; r3: cumple \\
icmp.seq & incremento 1 por paquete & r2: cumple; r3: cumple \\
icmp.checksum & variable en la captura & r2: no exacto; r3: no exacto \\
\bottomrule
\end{longtable}
\endgroup
Payload: [16] bytes. Hexadecimal inicial (offset 0):

\begin{quote}\ttfamily\footnotesize\raggedright 77fd 0000 0000 0000 0000 0000 0000 0000\end{quote}
La suma de payload\_\allowbreak{}word0\_\allowbreak{}BE, icmp.checksum, icmp.ident, icmp.seq módulo 65535 se mantuvo constante. La base se aprende del primer paquete de cada ejecución; no se presume universal. Es compatible con compensación de checksum, sin probar la intención del emisor.

Predicción concreta en ronda 2, campo icmp.seq: la primera solicitud tiene valor 0. Con la regla aprendida y esa base cuando corresponde, se predicen las solicitudes 2 a 4 antes de contrastarlas con sus valores observados.

\begingroup
\fontsize{9}{11}\selectfont
\addtocounter{table}{-1}
\begin{longtable}{p{3.2cm}p{5.1cm}p{5.1cm}}
\toprule
\textbf{Solicitud / trama} & \textbf{Predicción} & \textbf{Observación} \\ \midrule
\endfirsthead
\toprule
\textbf{Solicitud / trama} & \textbf{Predicción} & \textbf{Observación} \\ \midrule
\endhead
2 / 3 & 1 & 1 \\
3 / 5 & 2 & 2 \\
4 / 7 & 3 & 3 \\
\bottomrule
\end{longtable}
\endgroup
Alcance: 44 solicitudes en el ejemplo inicial. Una verificación de campo o byte no equivale a una medición independiente; los campos variables conservan sus restricciones y no reciben valores inventados.

\Needspace{9\baselineskip}
\subsubsection{scamper / tracelb / udp-sport}
Ejemplo: Valparaíso. Entrenamiento: scamper\_\allowbreak{}20261026T204908\_\allowbreak{}3c31fd8c; contraste: rondas 2 y 3. Las reglas completas de los tres destinos están en el archivo de predicciones enlazado en el anexo.

\begingroup
\fontsize{9}{11}\selectfont
\addtocounter{table}{-1}
\begin{longtable}{p{3.1cm}p{5.8cm}p{5.1cm}}
\toprule
\textbf{Campo} & \textbf{Regla en ronda 1} & \textbf{Contraste posterior} \\ \midrule
\endfirsthead
\toprule
\textbf{Campo} & \textbf{Regla en ronda 1} & \textbf{Contraste posterior} \\ \midrule
\endhead
ip.ttl & variable en la captura & r2: no exacto; r3: no exacto \\
ip.id & incremento 1 por paquete & r2: cumple; r3: cumple \\
ip.len & constante 44 & r2: cumple; r3: cumple \\
ip.proto & constante 17 & r2: cumple; r3: cumple \\
udp.srcport & variable en la captura & r2: no exacto; r3: no exacto \\
udp.dstport & constante 33435 & r2: cumple; r3: cumple \\
udp.length & constante 24 & r2: cumple; r3: cumple \\
udp.checksum & incremento 1 por paquete & r2: cumple; r3: cumple \\
\bottomrule
\end{longtable}
\endgroup
Payload: [16] bytes. Hexadecimal inicial (offset 0):

\begin{quote}\ttfamily\footnotesize\raggedright 66b1 0000 0000 0000 0000 0000 0000 0000\end{quote}
La suma de payload\_\allowbreak{}word0\_\allowbreak{}BE, udp.checksum, udp.srcport, udp.dstport módulo 65535 se mantuvo constante. La base se aprende del primer paquete de cada ejecución; no se presume universal. Es compatible con compensación de checksum, sin probar la intención del emisor.

Predicción concreta en ronda 2, campo ip.id: la primera solicitud tiene valor 0x0001. Con la regla aprendida y esa base cuando corresponde, se predicen las solicitudes 2 a 4 antes de contrastarlas con sus valores observados.

\begingroup
\fontsize{9}{11}\selectfont
\addtocounter{table}{-1}
\begin{longtable}{p{3.2cm}p{5.1cm}p{5.1cm}}
\toprule
\textbf{Solicitud / trama} & \textbf{Predicción} & \textbf{Observación} \\ \midrule
\endfirsthead
\toprule
\textbf{Solicitud / trama} & \textbf{Predicción} & \textbf{Observación} \\ \midrule
\endhead
2 / 9 & 2 & 0x0002 \\
3 / 24 & 3 & 0x0003 \\
4 / 29 & 4 & 0x0004 \\
\bottomrule
\end{longtable}
\endgroup
Alcance: 48 solicitudes en el ejemplo inicial. Una verificación de campo o byte no equivale a una medición independiente; los campos variables conservan sus restricciones y no reciben valores inventados.

\Needspace{9\baselineskip}
\subsubsection{scamper / tracelb / tcp-sport}
Ejemplo: Valparaíso. Entrenamiento: scamper\_\allowbreak{}20261026T205303\_\allowbreak{}343973b5; contraste: rondas 2 y 3. Las reglas completas de los tres destinos están en el archivo de predicciones enlazado en el anexo.

\begingroup
\fontsize{9}{11}\selectfont
\addtocounter{table}{-1}
\begin{longtable}{p{3.1cm}p{5.8cm}p{5.1cm}}
\toprule
\textbf{Campo} & \textbf{Regla en ronda 1} & \textbf{Contraste posterior} \\ \midrule
\endfirsthead
\toprule
\textbf{Campo} & \textbf{Regla en ronda 1} & \textbf{Contraste posterior} \\ \midrule
\endhead
ip.ttl & variable en la captura & r2: no exacto; r3: no exacto \\
ip.id & incremento 1 por paquete & r2: cumple; r3: cumple \\
ip.len & constante 40 & r2: cumple; r3: cumple \\
ip.proto & constante 6 & r2: cumple; r3: cumple \\
tcp.srcport & variable en la captura & r2: no exacto; r3: no exacto \\
tcp.dstport & constante 80 & r2: cumple; r3: cumple \\
tcp.flags & constante 0x0002 & r2: cumple; r3: cumple \\
tcp.seq\_\allowbreak{}raw & constante 0 & r2: cumple; r3: cumple \\
tcp.ack\_\allowbreak{}raw & constante 0 & r2: cumple; r3: cumple \\
tcp.len & constante 0 & r2: cumple; r3: cumple \\
tcp.checksum & variable en la captura & r2: no exacto; r3: no exacto \\
\bottomrule
\end{longtable}
\endgroup
Payload: [0] bytes. No hay datos de aplicación en las sondas TCP observadas.

Predicción concreta en ronda 2, campo ip.id: la primera solicitud tiene valor 0x0001. Con la regla aprendida y esa base cuando corresponde, se predicen las solicitudes 2 a 4 antes de contrastarlas con sus valores observados.

\begingroup
\fontsize{9}{11}\selectfont
\addtocounter{table}{-1}
\begin{longtable}{p{3.2cm}p{5.1cm}p{5.1cm}}
\toprule
\textbf{Solicitud / trama} & \textbf{Predicción} & \textbf{Observación} \\ \midrule
\endfirsthead
\toprule
\textbf{Solicitud / trama} & \textbf{Predicción} & \textbf{Observación} \\ \midrule
\endhead
2 / 3 & 2 & 0x0002 \\
3 / 7 & 3 & 0x0003 \\
4 / 13 & 4 & 0x0004 \\
\bottomrule
\end{longtable}
\endgroup
Alcance: 38 solicitudes en el ejemplo inicial. Una verificación de campo o byte no equivale a una medición independiente; los campos variables conservan sus restricciones y no reciben valores inventados.

\Needspace{9\baselineskip}
\subsubsection{scamper / tracelb / tcp-ack-sport}
Ejemplo: Valparaíso. Entrenamiento: scamper\_\allowbreak{}20261026T205617\_\allowbreak{}88e8814e; contraste: rondas 2 y 3. Las reglas completas de los tres destinos están en el archivo de predicciones enlazado en el anexo.

\begingroup
\fontsize{9}{11}\selectfont
\addtocounter{table}{-1}
\begin{longtable}{p{3.1cm}p{5.8cm}p{5.1cm}}
\toprule
\textbf{Campo} & \textbf{Regla en ronda 1} & \textbf{Contraste posterior} \\ \midrule
\endfirsthead
\toprule
\textbf{Campo} & \textbf{Regla en ronda 1} & \textbf{Contraste posterior} \\ \midrule
\endhead
ip.ttl & variable en la captura & r2: no exacto; r3: no exacto \\
ip.id & incremento 1 por paquete & r2: cumple; r3: cumple \\
ip.len & constante 40 & r2: cumple; r3: cumple \\
ip.proto & constante 6 & r2: cumple; r3: cumple \\
tcp.srcport & variable en la captura & r2: no exacto; r3: no exacto \\
tcp.dstport & constante 80 & r2: cumple; r3: cumple \\
tcp.flags & constante 0x0010 & r2: cumple; r3: cumple \\
tcp.seq\_\allowbreak{}raw & constante 0 & r2: cumple; r3: cumple \\
tcp.ack\_\allowbreak{}raw & constante 0 & r2: cumple; r3: cumple \\
tcp.len & constante 0 & r2: cumple; r3: cumple \\
tcp.checksum & variable en la captura & r2: no exacto; r3: no exacto \\
\bottomrule
\end{longtable}
\endgroup
Payload: [0] bytes. No hay datos de aplicación en las sondas TCP observadas.

Predicción concreta en ronda 2, campo ip.id: la primera solicitud tiene valor 0x0001. Con la regla aprendida y esa base cuando corresponde, se predicen las solicitudes 2 a 4 antes de contrastarlas con sus valores observados.

\begingroup
\fontsize{9}{11}\selectfont
\addtocounter{table}{-1}
\begin{longtable}{p{3.2cm}p{5.1cm}p{5.1cm}}
\toprule
\textbf{Solicitud / trama} & \textbf{Predicción} & \textbf{Observación} \\ \midrule
\endfirsthead
\toprule
\textbf{Solicitud / trama} & \textbf{Predicción} & \textbf{Observación} \\ \midrule
\endhead
2 / 3 & 2 & 0x0002 \\
3 / 5 & 3 & 0x0003 \\
4 / 6 & 4 & 0x0004 \\
\bottomrule
\end{longtable}
\endgroup
Alcance: 27 solicitudes en el ejemplo inicial. Una verificación de campo o byte no equivale a una medición independiente; los campos variables conservan sus restricciones y no reciben valores inventados.

